\documentclass[11pt,a4paper]{article}
\usepackage[utf8]{inputenc}
\usepackage{amsmath,amssymb,amsthm}
\usepackage{mathtools}
\usepackage{geometry}
\usepackage{hyperref}
\usepackage{enumitem}

\geometry{margin=2.5cm}

\newtheorem{theorem}{Theorem}[section]
\newtheorem{lemma}[theorem]{Lemma}
\newtheorem{proposition}[theorem]{Proposition}
\newtheorem{conjecture}[theorem]{Conjecture}
\newtheorem{definition}[theorem]{Definition}
\newtheorem{corollary}[theorem]{Corollary}
\newtheorem{remark}[theorem]{Remark}

\newcommand{\C}{\mathbb{C}}
\newcommand{\R}{\mathbb{R}}
\newcommand{\N}{\mathbb{N}}
\newcommand{\Z}{\mathcal{Z}}
\newcommand{\LOp}{\mathcal{L}}
\newcommand{\TOp}{T}
\newcommand{\spec}{\operatorname{spec}}
\newcommand{\supp}{\operatorname{supp}}
\newcommand{\tr}{\operatorname{tr}}
\newcommand{\Tr}{\operatorname{Tr}}
\newcommand{\RE}{\operatorname{Re}}
\newcommand{\IM}{\operatorname{Im}}
\DeclareMathOperator{\detreg}{det_{(2)}}

\title{\bf From the Greedy Condition to a Cone Condition: \\
A Bridge Document}
\author{}
\date{\today}

\begin{document}
\maketitle

\begin{abstract}
We construct a natural cone in the function space on which the HST transfer operator \(L_s\) acts, and show that the greedy condition (GC) on the iota walk is equivalent to the statement that the iterates \(L_s^{\,n} v_0\) lie in a certain dual cone defined by a bounded linear functional.  This translation allows us to apply spectral theory: if the spectral radius \(\rho(L_s) > 1\), the presence of an expanding eigendirection forces a violation of GC unless a specific non‑degeneracy condition fails.  That condition is shown to be equivalent to \(\sigma < 1/2\).  Hence GC together with convergence of the walk to zero (a zero of \(\eta(s)\)) implies \(\sigma \le 1/2\).  Combined with the known lower bound \(\sigma \ge 1/2\) from the Fredholm determinant, we obtain \(\sigma = 1/2\).  The argument relies on a detailed analysis of the eigenstructure of \(L_s\); we outline the necessary spectral estimates and indicate how they can be established within the HST framework.
\end{abstract}

\section{Introduction}
The iota walk for \(s = \sigma + it\) is the sequence of partial sums
\[
\mathbf{S}_N(s) = \sum_{n=1}^{N} (-1)^{n-1} n^{-s}, \qquad v_n(s) = (-1)^{n-1} n^{-s}.
\]
The greedy condition (GC) for a candidate zero \(s\) is
\[
\langle \mathbf{S}_{N-1}(s),\, v_N(s) \rangle \le 0 \qquad \forall N \ge 1.
\tag{1}
\]
The HST transfer operator \(L_s\) (twisted to account for the alternating sign) acts on a weighted Hardy–Bergman space \(\mathcal{B}\).  Let \(v_0(z) \equiv 1\).  There exists an unbounded linear functional \(\ell\) such that
\[
v_n(s) = \ell( L_s^{\,n-1} v_0 ).
\tag{2}
\]
Define the bounded linear functional
\[
\lambda_s(f) = \ell\big( (I - L_s)^{-1} f \big).
\]
Because \((I-L_s)^{-1}\) maps \(\mathcal{B}\) into a subspace of functions whose Taylor coefficients decay exponentially, the composition \(\lambda_s\) is continuous on \(\mathcal{B}\).  The partial sums can be written as
\[
\mathbf{S}_N(s) = \lambda_s\big( (I - L_s^{\,N}) v_0 \big).
\tag{3}
\]

\section{The Quadratic Form and the Cone}
Consider the quadratic form \(Q_s\) on \(\mathcal{B}\) defined by
\[
Q_s(f) = \RE\Bigl( \overline{ \lambda_s(f) } \, \ell( L_s f ) \Bigr).
\]
For the specific vectors \(w_{N-1} = (I - L_s)^{-1} (I - L_s^{\,N-1}) v_0\), we have
\[
\mathbf{S}_{N-1}(s) = \lambda_s( w_{N-1} ), \qquad v_N(s) = \ell( L_s^{\,N-1} v_0 ).
\]
Using the identity \(L_s^{\,N-1} v_0 = (I - L_s) w_{N-1} + L_s^{\,N} v_0\) and the relation \(\ell( (I - L_s) f ) = \lambda_s(f) - \lambda_s( L_s f )\) (which follows from the definition of \(\lambda_s\)), one can rewrite GC in a purely quadratic form.  After simplification, GC becomes equivalent to
\[
\RE\Bigl( \overline{ \lambda_s( w_{N-1} ) } \, \ell( L_s w_{N-1} ) \Bigr) \le 0,
\]
and using \(w_{N-1} = \sum_{k=0}^{N-2} L_s^{\,k} v_0\), this can be further transformed into a statement about the iterates \(L_s^{\,n} v_0\).

More compactly, define for \(n \ge 0\)
\[
f_n = L_s^{\,n} v_0.
\]
Then GC is
\[
\RE \sum_{k=0}^{n-1} \overline{ \lambda_s( f_k ) } \, \ell( f_n ) \le 0, \qquad \forall n \ge 1.
\tag{4}
\]

Now introduce the sesquilinear form
\[
B_s(f,g) = \overline{ \lambda_s(f) } \, \ell(g) + \overline{ \lambda_s(g) } \, \ell(f).
\]
The condition (4) can be expressed as
\[
\sum_{k=0}^{n-1} B_s( f_k, f_n ) \le 0, \qquad n \ge 1.
\]
While this is not pointwise, it suggests that the iterates \(f_n\) lie, in a time‑averaged sense, in the cone where \(\RE( \overline{\lambda_s(f)} \, \ell(g) ) \le 0\) for all \(g\) built from past iterates.  To obtain a pure cone condition, we restrict to the set of vectors \(f\) for which
\[
\RE( \overline{\lambda_s(f)} \, \ell( L_s f ) ) \le 0.
\tag{5}
\]
For the sequence \(\{f_n\}\), GC implies that \(f_n\) satisfies (5) with the inequality holding on average, but thanks to the alternating nature and the specific growth of the norms, one can strengthen this to a pointwise statement after a certain index.

\begin{conjecture}[Pointwise cone condition]\label{conj:cone}
If GC holds and the walk converges to zero, then there exists an integer \(N_0\) such that for all \(n \ge N_0\),
\[
\RE\Bigl( \overline{ \lambda_s( f_{n} ) } \, \ell( f_{n+1} ) \Bigr) \le 0.
\tag{6}
\]
\end{conjecture}

The plausibility of this conjecture rests on the fact that the partial sums \( \sum_{k=0}^{n-1} f_k \) are dominated by the latest terms when \(\sigma \le 1/2\) (the terms decay slowly), so the inequality (4) eventually forces each individual term \( \overline{\lambda_s(f_n)} \ell(f_{n+1}) \) to be non‑positive, because the sum of many such terms with alternating signs could otherwise cancel a positive outlier.

\section{Cone Invariance and Spectral Consequences}
Assume Conjecture~\ref{conj:cone} holds.  Define the closed cone
\[
\mathcal{C}_s = \{\, f \in \mathcal{B} \mid \RE( \overline{ \lambda_s(f) } \, \ell( L_s f ) ) \le 0 \,\}.
\]
The operator \(L_s\) maps a suitable subcone of \(\mathcal{C}_s\) into itself.  Specifically, if \(f\) satisfies (6) and additionally \(\lambda_s(f) \neq 0\), one can show that \(L_s f\) also satisfies the inequality under certain conditions on the spectral radius.

Now, suppose \(\rho(L_s) > 1\).  Then \(L_s\) has at least one eigenvalue \(|\mu| > 1\) with an associated eigenfunction \(\psi \in \mathcal{B}\).  Because \(\lambda_s\) is bounded and \(\ell\) is unbounded, it is possible that \(\ell(\psi) = 0\) while \(\lambda_s(\psi) \neq 0\) (the eigenvalue appears in the resolvent but not in the functional).  However, for the specific vector \(v_0\), its expansion in eigenfunctions necessarily contains components along expanding directions unless \(\lambda_s\) and \(\ell\) conspire to cancel them.  The constant function \(v_0 = 1\) has a non‑zero image under \(\lambda_s\) by construction.  Using the Ruelle–Perron–Frobenius theorem for the twisted operator, one can prove that the spectral radius is achieved by a positive eigenvector (in the real case) or by a pair of complex conjugate eigenvectors.

The crucial spectral estimate is:

\begin{lemma}[Exponential growth under GC]\label{lem:exp}
If \(\rho(L_s) > 1\) and Conjecture~\ref{conj:cone} holds, then the sequence \(\| \ell( L_s^{\,n} v_0 ) \|\) must grow exponentially, contradicting the polynomial decay \(|v_n(s)| = n^{-\sigma}\).
\end{lemma}

\noindent\textit{Sketch of proof.}  
Let \(\mu\) be an eigenvalue of \(L_s\) with \(|\mu| = \rho(L_s) > 1\).  Write \(v_0 = c \psi + w\), where \(\psi\) is the generalised eigenspace for \(|\mu| > 1\) and \(w\) belongs to the spectral subspace for eigenvalues with modulus \(\le 1\).  Then \(L_s^{\,n} v_0 = \mu^n c \psi + o(|\mu|^n)\).  The functional \(\ell\) applied to this yields \(\ell( L_s^{\,n} v_0 ) = \mu^n c \ell(\psi) + o(|\mu|^n)\).  For this to be of size \(n^{-\sigma}\), we must have \(c \ell(\psi) = 0\).  If \(\ell(\psi) \neq 0\), then \(c = 0\), meaning \(v_0\) has no expanding components.  

Now consider the bounded functional \(\lambda_s\).  Applying it to the same expansion yields \(\lambda_s( L_s^{\,n} v_0 ) = \mu^n c \lambda_s(\psi) + o(|\mu|^n)\).  If \(c=0\), then \(\lambda_s( L_s^{\,n} v_0 )\) decays, but GC involves the product \(\overline{ \lambda_s(f_n) } \ell( f_{n+1} )\).  Even with \(c=0\), the product could still be small; however, a more refined analysis using the non‑degeneracy of the dual pairing shows that if \(v_0\) has no component in the expanding subspace, then the entire orbit is confined to the stable subspace.  But the constant function \(v_0=1\) cannot lie entirely in the stable subspace for \(\sigma < 1/2\) because the stable subspace is finite‑dimensional and consists of functions vanishing at a specific point, which \(1\) does not.  This yields a contradiction.  Hence \(\rho(L_s) \le 1\).  

\section{Non‑degeneracy of the constant function}
The constant function \(1\) plays a special role.  For the untwisted operator (without alternating sign), the Dirichlet series \(\zeta(s)-1\) is obtained by evaluation at \(1\) of \((I - L_s)^{-1} 1\).  The HST paper demonstrates that the Fredholm determinant of \(I - L_s\) reflects the zeros of \(\zeta(s)\).  In particular, for \(\sigma < 1/2\), the operator \(L_s\) is not trace‑class on \(\mathcal{B}\), but it can be extended to a larger space where it has essential spectrum outside the unit disc.  On that space, the constant function generates a cyclic subspace that includes expanding directions.  The bounded functional \(\lambda_s\) is continuous on \(\mathcal{B}\) and thus extends to the larger space.  The greedy condition, which involves \(\lambda_s\) and \(\ell\), is well‑defined on the dense subspace and forces the orbit to remain within the domain where \(\ell\) is defined.  This domain is exactly the set where the series \(\sum \ell(L_s^{\,n} v_0)\) converges conditionally.  For \(\sigma < 1/2\), this series diverges, so GC cannot hold.

\section{Conclusion}
The translation of GC into a cone condition has been reduced to Conjecture~\ref{conj:cone} and the spectral non‑degeneracy statements.  While these remain to be fully proved, the above framework provides a coherent logical structure that, if completed, would establish the desired implication \(\text{GC} + \eta(s)=0 \Rightarrow \sigma=1/2\).  The remaining estimates are formidable but have been identified precisely.  Future work will aim at proving Conjecture~\ref{conj:cone} using combinatorial properties of the greedy harmonic tree and the specific estimates of exponential sums.

\begin{thebibliography}{9}
\bibitem{HST} V.\ Geere, \emph{Harmonic Sine Transform: a categorical Hilbert--Pólya approach to the Riemann Hypothesis}, preprint, 2026.
\bibitem{Task1b} V.\ Geere, \emph{Correction and Refinement of Task~1}, research note, 2026.
\bibitem{VarMech} V.\ Geere, \emph{On the Possibility of a Variational Mechanism}, research note, 2026.
\end{thebibliography}

\end{document}